\section{Define the Job, Decision, and Operating Context Before Decomposing the Tool}

\textbf{Why this matters:} A system can complete its assigned task perfectly and still fail the reason it was built, because the task, the output, and the actual job being served are not the same thing.

\textbf{Practical implication:} Before assessing any tool technically, write down what decision it's meant to inform and who is actually authorized to rely on its output. Most disputes about whether a tool "works" trace back to skipping this step.

A product is suitable, or not, only for a particular job performed under particular conditions. The relevant unit of analysis is therefore not the product in the abstract, and not even the task considered in isolation, but the complete relationship between an intended outcome, the work performed to reach it, the people who rely on that work, and the consequences that follow.

A task description alone is insufficient, and this is worth taking seriously rather than treating as a pedantic distinction. "Summarize these contracts," "research this legal issue," or "draft a response" describes an activity, but not why the activity is being performed or what would make it useful. A system can perform the stated activity accurately and still fail the underlying job: it can produce a correct summary that omits the issues material to the pending decision, sound research that arrives after the decision has already been made, or a well-written response that implies an authority the sender doesn't actually possess. None of that shows up as an error under Section 9's diagnostic ladder --- the output can pass every one of those five axes and still be \textbf{a successful production of the wrong operational result.}

Five related concepts are worth distinguishing explicitly, as definitions this report adopts rather than as claims requiring evidence:

\begin{itemize}
\item \textbf{The job} is the progress or outcome actually being sought.
\item \textbf{The decision} is the judgment, choice, or action the work is intended to inform.
\item \textbf{The task} is the activity assigned to the human or the system.
\item \textbf{The output} is the artifact, recommendation, classification, or state change produced.
\item \textbf{The effect} is what happens after the output is accepted and used.
\end{itemize}

These shouldn't be treated as interchangeable, and collapsing them is the single most common way a job-level failure gets misdiagnosed as a technical one. The output is not the outcome, and completing the task is not proof the job has been accomplished. A workflow that optimizes the production of deliverables can look highly successful --- high volume, high acceptance rate, positive user feedback --- while making little or no actual contribution to the decision or outcome those deliverables were requested to inform. Consider a system asked to summarize a hundred contracts (\textbf{task}) and produce a table of clauses and extracted terms (\textbf{output}), where the real purpose is to let a deal team identify which contracts require escalation before a transaction closes (\textbf{decision}), so the team can change the transaction plan early enough for that to matter (\textbf{job}), resulting in issues actually being escalated, negotiated, priced, or knowingly accepted (\textbf{effect}). A system can nail the task and the output --- accurate, well-formatted, on time by some measure --- while failing the job outright, because it didn't weight transaction-critical provisions correctly, finished too late to change anything, applied the wrong materiality threshold, or presented results in a form the actual decision-makers couldn't act on. That failure is neither a factual error nor a context error in Section 9's terms. It's a correctly executed task in service of the wrong job.

A useful job definition begins with the triggering circumstance and the progress actually sought, in a single sentence:

\begin{quote}
When \textbf{[a defined circumstance arises]}, \textbf{[a defined actor]} needs to \textbf{[make progress or reach a decision]}, using \textbf{[specified information and expertise]}, so that \textbf{[the intended operational outcome]} can occur.

\end{quote}

That single sentence establishes intent, but intent alone doesn't carry everything the rest of this framework needs --- a second artifact is required to establish operational meaning: a short use-and-reliance description. At minimum, it should identify: the requester, operator, reviewer, decision-maker, affected parties, and the named owner of residual risk (these are frequently different people, and treating them as one undifferentiated "user" is where a lot of the trouble in the rest of this list originates); the decision or action the output will actually inform; the required information, permissions, expertise, and temporal freshness the job depends on; the form the output must take to be genuinely useful to the person relying on it; the degree of reliance that may properly be placed on it; the review or approval required before it travels any further; the validated operating envelope and its known exclusions, in the sense already established in Section 2; the circumstances requiring clarification, narrowing, abstention, or escalation; the downstream systems, records, or people that will receive the result; and the measures by which success and failure will actually be judged.

The distinction between the operator and the beneficiary deserves particular attention, because it's an easy one to collapse without noticing. The person interacting with the system may not be the person whose objective the workflow actually exists to serve. A junior lawyer may operate a contract-review system, a partner may rely on its synthesis, a client may act on the partner's advice, and third parties may be affected by the resulting transaction --- four distinct roles, only one of which is "the user" in the ordinary sense. Defining the job solely from the operator's perspective routinely obscures the more consequential forms of reliance sitting further downstream, where the actual stakes usually live.

The same collapsing risk applies to authority specifically. A person may be fully permitted to request or review an output without being authorized to accept the risk it identifies, communicate it externally, or cause an action to occur on its basis. The job definition should state not merely who uses the system, but who is actually authorized to decide, approve, communicate, and act --- since a system that quietly lets an unauthorized person do any of those four things has failed regardless of how accurate its output was.

The current human workflow should be recorded before the AI-enabled alternative is designed, not after. This does two things at once: it establishes the strong human comparator Section 2 requires, and it surfaces what the existing process is actually accomplishing beyond its stated purpose. Apparent inefficiencies in a human process are frequently performing a hidden function --- creating time for reflection, forcing a second person to actually inspect the record, preserving a legible chain of responsibility, or surfacing disagreement that a faster process would silently suppress instead of resolving. The goal isn't to reproduce every existing step for its own sake, but to understand which outcomes and controls those steps were actually providing before removing or replacing them, so that a faster workflow doesn't quietly discard a control nobody noticed was doing real work.

Success should be defined at the level of the job, not only at the level of the output. Depending on the job, success might mean fewer material omissions, earlier identification of risk, more consistent treatment of equivalent cases, lower rework, faster completion, improved decision quality, or reduced demand on scarce expert attention. Output completion, user adoption, volume processed, and reviewer acceptance rate are all useful operational signals, but none of them, alone, establishes that the underlying job is actually being done well --- a system can score highly on every one of them while the job itself goes unserved.

Failure, correspondingly, has to include more than incorrect content. A workflow can fail because the answer arrived late, because it was presented to an audience that couldn't act on it, because it created more verification work downstream than it saved upstream, because it exceeded the user's actual authority, or because it caused someone to place more reliance on the result than the evidence available actually justified. These are job-level failures even where the output itself, examined on its own terms, is entirely factually sound --- which is exactly why Section 9's output-centered ladder, useful as it is, cannot catch them on its own.

The boundaries of the job need to be explicit, not assumed. The assessment should identify where the AI-enabled workflow begins and ends --- whether it retrieves, extracts, analyzes, recommends, drafts, approves, communicates, or acts, and which of those functions remain firmly outside its authority regardless of what it's technically capable of doing. It should also name the adjacent uses that have specifically \emph{not} been approved. A system validated as a research assistant doesn't become a legal-advice system merely because a user copies its output into an advice note; a system validated for first-pass drafting doesn't become a communications agent merely because it happens to have technical access to an email tool.

This creates a distinct and easy-to-miss risk worth naming directly: \textbf{job drift.} A system validated for one job gets gradually used for adjacent jobs, not through any deliberate decision, but because its outputs look useful and the technical interface happens to permit the expansion. A research assistant becomes an advice generator; a drafting tool becomes a final-communication tool; an issue-spotting system becomes a risk-rating system; an internal summarizer becomes a client-facing reporting tool. The underlying software, and often the underlying task, can look essentially unchanged throughout --- but the reliance placed on it, the authority it's now operating under, the strength of verification it actually has, and the consequences of its errors have all shifted, frequently without anyone deciding that they should. Job drift is, in effect, a form of distribution shift at the workflow level --- the same underlying mechanism Section 4 flags for a checker whose training distribution no longer matches its deployment distribution, and Section 7 flags for a model whose specification difficulty has quietly increased, now occurring in the surrounding organizational process rather than in the model itself. Worth building in as controls against it: a named, specifically approved job, stated narrowly enough to be checked against; explicit adjacent uses that are excluded rather than merely unaddressed; visible output labeling that signals what the system was and wasn't validated for; access or action restrictions that make unapproved uses actually harder, not just discouraged; monitoring for usage patterns that don't match the approved job; and a trigger for re-evaluation whenever reliance changes, even where the underlying software hasn't changed at all.

Finally, a defensible conclusion should make clear the job being approved, the conditions under which the evidence applies, the degree of reliance justified, the principal residual uncertainties, and the circumstances that would require reconsideration. That accommodates a genuine range of outcomes --- a narrow pilot, a staged rollout, continued experimentation short of reliance, an outright no --- without implying every discussion of an AI-enabled workflow must end in a governance label.

How much of the rest of this report's machinery any given job actually needs is itself a judgment call, and it's worth stating the principle behind that judgment explicitly rather than leaving it implicit: \textbf{the depth of analysis should scale with the consequence, irreversibility, scale, sensitivity, and opacity of the workflow --- not simply with the fact that AI is involved.} A handful of factors reliably justify deeper analysis, and their absence justifies a lighter touch: sensitive or privileged information; external or client reliance; irreversible action; weak verification (Section 7); a broad action space (Section 6); high-volume reuse; persistent organizational memory (Section 14); and a large blast radius (Section 13). A low-stakes internal drafting aid touching none of these doesn't need the full apparatus in Sections 11 through 14 any more than it needs a formal use-and-reliance description; a workflow touching several of them at once does. The point of naming the factors is to let a reader scale their own effort deliberately, rather than either skipping the depth this report argues for or applying all of it uniformly regardless of what's actually at stake.

The operative question, then, is not merely "what task will the system perform," but:

\begin{quote}
What progress is being sought, what decision or action will the system's work inform, what information and authority does that require, how will the result actually be relied upon, and what would demonstrate that the complete workflow does the underlying job better than the competent human process it's intended to improve?

\end{quote}

A stronger treatment of the job gives every later section in this report a more precise role to play, rather than leaving each one to reconstruct "the job" for itself: Section 2 establishes the framework's overall standard of success; this section defines the particular job and use contract that standard gets applied to; Section 4 decomposes how a system technically attempts that job; Section 7 assesses how difficult that specific job is to specify and verify; Section 9 diagnoses failures in the individual outputs the job produces along the way; Section 11 identifies the unacceptable losses specifically associated with that job; Section 12 traces how a faulty state actually propagates through that job's workflow; Section 13 tests whether the evidence gathered actually supports the reliance the job requires; and Section 14 determines whether the job keeps being performed reliably once it's actually running in operation. None of those later sections can do that work well against an undefined job.
